\begin{titlepage}
\centering\vspace*{24mm}
{\Huge\bfseries Agentic State\par}
\vspace{5mm}
{\LARGE Context, Memory, and Control\\in Adaptive AI\par}
\vspace{10mm}
{\large A Technical Monograph on Conditioning,\\Tool Execution, Verification, and Self-Improvement\par}
\vfill
{\large Technical Monograph\par}
\vspace{8mm}
{\normalsize Literature cutoff: 29 October 2026\par}
{\normalsize Latest source retrieval: 29 October 2026\par}
\end{titlepage}

\chapter*{Abstract}
This monograph analyzes agentic computation as a versioned, partially observed control system. Across 34 chapters, it connects model architecture and learning to context construction, session and persistent memory, procedural skills, tool interfaces, orchestration, evaluation, security, alignment, self-improvement, and deployment. The unit of analysis is the coupled system: a model checkpoint alone does not identify the information available to a decision, the operation permitted to execute, or the evidence supporting its outcome.

The central synthesis separates proposal-law preservation, authority nonexpansion, and continuity of operation identities and effect status. These contracts distinguish a transformation that preserves state from a system that makes correct decisions. Conditional-information and change-of-measure arguments identify the assumptions needed for bounded decision and trajectory comparisons. Training likelihoods must distinguish the actual decoder and behavior-time inputs from a reconstructed transcript. Trusted authorization and effect state are distinguished from model-visible evidence; checkpoint restoration must not revive revoked grants, and a missing receipt does not prove that an external effect is absent. These distinctions connect context--memory--policy consistency to retrieval, cache validity, skill migration, and service-supported reconciliation.

Established results are distinguished from original synthesis, analytical deductions, implementation contracts, counterexamples, and proposed experiments. Author-reported findings retain their source and population limits; source inspection is not experimental reproduction. Model and harness comparisons specify adaptation, evaluator access, baselines, and complete resource boundaries. Resource analyses account for computation, parameter and optimizer memory, KV growth, bandwidth, latency, throughput, and cost under declared hardware and workload assumptions. No new experimental measurements are reported. Finite benchmark gains and improvement trajectories do not establish unrestricted AGI, arbitrary interface robustness, or indefinite recursive improvement.

\chapter*{Notation and evidence}
The default symbols are $H_t$ for recorded history, $C_t$ for working context, $M_t$ for persistent memory, $E_t$ for environment state, $S_t$ for a specified system state, $\theta$ for model-policy parameters, and $\delta$ for deployment decoding configuration. The compatibility manifest $\nu_t$ is distinct from the current authority epoch $\alpha_t$; $\lambda$ denotes an operation-effect ledger where specified. Local definitions take precedence where a chapter explicitly introduces another conventional meaning. A deployed proposal law includes its declared decoding, stopping, and parsing semantics; it need not equal a raw token likelihood or the law of authorized dispatches. Units and workload assumptions accompany resource calculations. A calculated example is not a measured benchmark result.

The literature cutoff and latest source-retrieval date are 29 October 2026 for the complete monograph. Foundational references support mathematics and mechanisms whose assumptions remain applicable; frontier claims are tied to identified papers, releases, and inspected revisions. An identifier, retrieval date, or repository inspection does not establish that a claim has been independently reproduced or remains valid under a different configuration.

Four evidence classes are distinguished: author-reported results, independently reproduced results, analytical deductions, and proposed experiments. This edition claims no independent experimental reproduction. Bibliographic identity checks, implementation inspection, and experimental reproduction are different evidential acts. Missing implementation details and unsupported inferences are identified rather than supplied by assumption. Proposed tests and mathematical bounds do not certify properties beyond their stated information, sampling, and enforcement premises.

\chapter*{Acknowledgments and computational assistance}
OpenAI Codex agents assisted source analysis, technical drafting, mathematical and editorial checks, and LaTeX preparation throughout Chapters 1--34 and the front matter. Explanatory illustrations were generated through OpenAI's built-in image-generation tool. The interface did not expose the image model identifier; no exact image-model version is asserted. The illustrations communicate mechanisms and contain no experimental measurements. These tools are not authors, and their use does not establish human expert endorsement. Responsibility for a submitted edition, its authorship, source verification, and final disclosures remains with its named human authors.
